\section{Discussion}

The preregistered question was whether breath-coupled visual feedback changed the relation between scheduled respiratory phase and audio--tactile SOA. The Condition $\times$ Phase $\times$ SOA interaction was not detected in the 40-person trial-level analysis. Treating SOA as categorical, excluding the documented partial session, allowing participant-cluster-robust uncertainty, and giving each participant one five-SOA slope contrast led to the same qualitative conclusion. The resulting confidence intervals still allow effects in either direction, so the finding should be described as a failure to detect the predicted modulation, not proof that respiration or visual feedback can have no PPS effect.

Audio--tactile responses were faster than tactile-only responses on average. The modest far-to-near RT gradient in Audio--Tactile trials was numerically similar to the gradient across nominal delay labels in Baseline trials (Fig.~\ref{study5:fig:delay_matched_facilitation}). Source recordings for 15 cleared participants found a 500-ms difference in actual trial-relative cue time between those trial types at the same label. Their label-based subtraction therefore cannot isolate or remove temporal expectancy, and the raw gradient cannot be read as a specific auditory-proximity effect. This concern is consistent with a four-experiment meta-analysis identifying expectancy, response probability, trial counts, and control conditions as consequential for interpreting near-versus-far tactile RT benefits \citep{holmes2020sounds}. The registered participant-level sigmoid gate failed in every trial-level cell, while four-parameter fits to only five SOA means were weakly identified and often reached parameter limits. The registration's approximate distance labels also disagree with the positions implied by its stated linear trajectory; replacing the axis left the fit gate unchanged but altered some fit diagnostics. Group-averaged curves can communicate the observed pattern, but a group midpoint is not an estimate of participant-specific PPS extent and cannot supply participant replicates for the planned boundary comparison. Smoothing or constraining curves after inspecting these data would change the estimand and should remain descriptive.

Several design and provenance features narrow the interpretation. Tactile-only Baseline RT varied with scheduled phase, undermining the assumed phase-invariant reference for Baseline subtraction. The primary test used raw Audio--Tactile RT for that reason. Respiratory phase labels identified the instructed half-cycle and were not uniformly verified against belt recordings at the trial level; a physiological phase effect cannot be isolated as cleanly as in an online-gated design. The documented cue-time mismatch may also alter the trial's position within the instructed breathing cycle. Source RT definitions and available original recordings differ across acquisition batches, and sound types changed between versions. Condition order was approximately balanced overall but related to dataset version. The hit-probability and trial-order checks did not reveal a different three-way pattern, but RTs on missed trials and 80 hit-labelled Audio--Tactile trials without a recorded RT remain unobserved. These factors do not explain away the primary null detection, yet they limit a precise claim about physiological PPS boundaries from the pooled curves.

The preregistration specified a Cauchy width for a Bayesian analysis without fully defining the mixed-model prior or random-effects comparison. The exploratory Bayes factor on participant slope contrasts favoured a point-null model under its stated assumptions, but it is not a substitute for the missing registered mixed-model Bayes factor. The scope of its low-sigmoid-fit participant exclusion was also ambiguous; we used it only to gate boundary-parameter inference. Eight pilot sessions were collected before registration as disclosed in the registration. The cleared pool contains 39 complete sessions and one admitted partial session, exceeding the registered target of 30 complete datasets; the reason should be documented before submission. Transparent reporting of these distinctions is preferable to presenting either a BIC approximation as the registered Bayes factor or the older 16-person, single-order result as the completed cohort.

A future boundary-focused study would benefit from more sampled SOAs around a candidate transition, harmonised response timing from retained recordings, trial-level physiological verification of respiratory phase, and a prespecified hierarchical curve model with parameter-recovery checks. For the present data, the most defensible result is the trial-level interaction test together with participant-weighted observed SOA means, explicit uncertainty, and the failed sigmoid-viability result.
